\documentclass[11pt]{article}
\usepackage[margin=1in]{geometry}
\usepackage[T1]{fontenc}
\usepackage[utf8]{inputenc}
\usepackage{lmodern}
\usepackage{setspace}
\usepackage{hyperref}
\usepackage{booktabs}
\usepackage{enumitem}
\usepackage{amsmath,amssymb}
\usepackage[numbers]{natbib}
\onehalfspacing

\title{From Robots to Legal Digital Personality: \\ Rethinking AI Status, Attribution, and Public Responsibility, 2019-2026}
\author{Researcher Agent English}
\date{April 3, 2026}

\begin{document}
\maketitle

\begin{abstract}
This article revisits a 2019 Swiss bachelor thesis on robot taxation and asks how its central legal intuition should be reformulated in 2026. The original thesis concluded that taxing robots as capital was generally inefficient, but suggested that taxation of machine-generated income would eventually require some form of electronic legal personality. In 2026, however, the legal landscape has shifted: European and Swiss AI regulation increasingly governs providers, deployers, and high-risk uses rather than granting legal personhood to artificial intelligence itself. This article argues that the thesis identified a real doctrinal problem but attached it to the wrong legal endpoint. The strongest 2026 framework is not full AI personhood, but a limited legal digital personality, understood as a functional attribution device for narrowly defined contexts where traceability, accountability, and selected fiscal consequences need to be organised with greater precision. Such a framework preserves human and corporate responsibility while recovering the most useful insight of the earlier robot-tax debate.
\end{abstract}

\noindent\textbf{Keywords:} artificial intelligence, robots, legal personality, digital personality, taxation, accountability, Switzerland

\section{Introduction}
The late-2010s debate on robot taxation was driven by a powerful intuition. If automation replaces labor, labor income may shrink, social-insurance financing may weaken, and fiscal systems built around employment may become less stable. The 2019 bachelor thesis at the center of this project entered that debate from a Swiss perspective. It examined whether taxing robots could preserve tax revenue and reduce pressure on labor, and concluded that such a tax was usually distortionary, administratively weak, and economically undesirable.

What makes the thesis still interesting in 2026 is not merely its anti-robot-tax conclusion. Its real importance lies in a more ambitious claim: once one stops thinking of robots as ordinary capital goods and starts attributing productivity, income, or quasi-autonomous action to them, one runs into the question of legal status. In the original text, this appeared as the possibility of an ``electronic personality.''

That move was intellectually ahead of its time, but it now needs reconstruction. Between 2019 and 2026, the legal center of gravity shifted away from physical robots and toward broader AI systems, including software, foundation models, and hybrid human-AI workflows. In parallel, major legal frameworks developed in Europe and Switzerland without embracing AI personhood. The result is a productive tension. The status question did not disappear, but the law answered it indirectly: by regulating actors, uses, and governance chains rather than AI entities as independent persons.

This article argues that the best way to extend the 2019 thesis is not to defend full AI personhood. It is to distinguish between full legal personhood, which remains doctrinally excessive, and a limited legal digital personality, which may be justified as a technical tool for attribution and responsibility in narrow domains. The argument proceeds in four steps. First, it recovers the logic of the original thesis. Second, it explains why the robot-centered frame is inadequate in 2026. Third, it shows why current law favors actor-based governance. Fourth, it develops a bounded concept of legal digital personality that does not dissolve human or corporate accountability.

\section{Literature Review}
The original thesis was built on the mature robot-tax literature of the late 2010s. Economists such as Acemoglu and Restrepo, Guerreiro, Rebelo and Teles, Thuemmel, and Gasteiger and Prettner explored automation, labor substitution, productivity, and the possible effects of taxing robots. Their shared lesson was not that automation posed no social problem, but that a robot tax often targeted the wrong object. In many models, taxing robots reduced efficiency, discouraged innovation, or merely shifted burdens elsewhere.

At the same time, legal scholars such as Xavier Oberson and Luciano Floridi pushed the debate into a different register. Their question was less about the immediate optimality of a tax and more about the legal conditions under which one could even speak coherently about machine-generated taxable capacity, duties, or accountability. That is where the notion of electronic or artificial personality appeared.

From a 2026 standpoint, both strands remain relevant but incomplete. The economics literature was right to distrust simple robot taxes, yet often treated legal status as secondary. The legal literature was right to identify a conceptual bottleneck, but sometimes overstated the attraction of personhood language. Since 2024, the dominant regulatory model has been neither laissez-faire nor AI personhood. It has instead been risk-based and actor-based. This makes the older literature newly valuable, because it allows us to see what problem survived once the proposed solution changed.

\section{From Robots to AI, 2019-2026}
The transition from 2019 to 2026 is not a simple update in terminology. It is a shift in the legal object under discussion. In 2019, the robot could still function as a visible proxy for automation: a machine, a production input, perhaps a quasi-worker in speculative argument. In 2026, AI is not adequately represented by that image. Relevant systems are often intangible, iterative, embedded in software infrastructure, licensed across borders, or integrated into ordinary business processes without clear physical boundaries.

This matters for three reasons. First, the fiscal base becomes harder to identify. A machine-specific levy is ill-suited to software, models, and cloud-based services. Second, the legal subject becomes harder to define. The more distributed the system, the less plausible it is to treat it as a self-standing unit comparable to a corporation or natural person. Third, the governance challenge changes. The key questions are no longer only about labor displacement and tax erosion, but also about transparency, human oversight, fundamental-rights risk, safety, traceability, and allocation of responsibility.

The original thesis partly anticipated this shift. It already suggested that if robots came to resemble workers in a legally meaningful way, then taxation would have to follow attributed income rather than mere existence. But its framework still relied too heavily on identifiable robots. The 2026 debate requires a more abstract and more careful category.

\section{What Legal Personality Means Here}
Legal personality is not a mystical status. It is a legal technique. To say that an entity has legal personality is to say that the law treats it as a point of attribution for rights, duties, capacities, liabilities, and legal acts. The law can do this for many reasons: practical convenience, asset partitioning, stable representation, responsibility allocation, or political value.

This immediately shows why the AI debate often becomes confused. People sometimes use ``legal personality'' in three different senses. The first is full personhood, comparable in legal effect to corporate personality. The second is symbolic personality, used rhetorically to express the social importance of AI. The third is functional attribution, which does not require the whole bundle of legal capacities but only some limited form of legal recognition.

For AI, these distinctions are decisive. Full personhood would imply a very large restructuring of private law, public law, tax law, and liability doctrine. It would require robust answers to questions of representation, patrimony, sanction, insolvency, and rights. Symbolic personhood adds little. Functional attribution, by contrast, is more promising. It asks whether the law needs a limited interface to connect a system's technical identity, deployment history, logs, and governance chain with existing human and corporate responsibility.

That narrower idea is what this article calls legal digital personality.

\section{Why Current Law Regulates Actors Rather than AI Persons}
The most important legal development since the original thesis is that the leading public-law frameworks for AI did not move toward AI personhood. They moved toward obligations imposed on providers, deployers, and other actors in the lifecycle of AI systems. This is not a conceptual accident; it reflects the structure of the problem.

Where the risk concerns safety, discrimination, opacity, manipulation, or rights violations, the law typically needs identifiable addressees who can document, explain, disclose, prevent, and remedy. Human beings, firms, authorities, and other organisations can perform those legal functions in a way present-day AI cannot. They can be supervised, fined, ordered to change conduct, and held answerable in existing institutional systems.

The legal consequence is clear. The law does not need AI personhood in order to regulate AI effectively. It needs robust governance chains. This is why risk-based regimes focus on classification, transparency, conformity, oversight, logging, and accountability. A system can be heavily regulated without becoming a legal person.

Still, that is not the end of the matter. Actor-based regulation leaves open a narrower question: whether certain AI-mediated contexts would benefit from a supplementary attribution device that is more precise than ordinary product or corporate categories but far less radical than personhood.

\section{Fiscal and Macroeconomic Implications of the Status Choice}
The original thesis remains strongest when it explains why robot taxation poorly fits the economic reality of automation. That argument becomes even stronger in the age of AI. If the relevant productive gains arise from software, data, model integration, and intangible scale effects, a tax keyed to robots as physical objects is conceptually weak and administratively fragile.

The status choice matters because it shapes where law looks for taxable capacity and responsibility. If AI is treated purely as an object within the assets of firms, then tax law will continue to focus on the conventional actors: the firm, the contractor, the employer, the platform. If AI were given full personality, tax law would have to solve an array of artificial attribution problems, including income measurement, asset ownership, and double taxation. The 2019 thesis already recognized this difficulty.

A limited legal digital personality offers a middle path. It does not create a new general taxpayer. It instead provides a way to stabilise the attribution chain around particular systems or uses. For example, it may support registration, event logging, identification of responsible operators, or ring-fenced accounting in narrow regulatory contexts. Fiscal consequences would remain primarily routed through human or corporate actors, but the legal mapping of system-generated effects could become more precise.

Macroeconomically, this matters because the central challenge is not to personify AI but to preserve institutional control over value creation and responsibility in an economy increasingly shaped by intangible, scalable, and partially autonomous systems. A well-designed legal digital personality could improve governance without pretending that AI itself is the final bearer of political or moral responsibility.

\section{A Proposal for Limited Legal Digital Personality}
A defensible legal digital personality should satisfy five constraints.

First, it must be \emph{functional rather than anthropomorphic}. The point is not to imitate human or corporate personhood. It is to solve practical attribution problems.

Second, it must be \emph{limited in scope}. It should be available only in contexts where ordinary categories leave a real gap. Not every AI system needs a special legal wrapper.

Third, it must be \emph{subordinate to human and corporate responsibility}. The creation of a digital personality should never become a shield that lets firms externalise liability to an incapable technological entity.

Fourth, it must be \emph{traceability-enhancing}. The category should improve recordkeeping, explainability of governance chains, and identification of competent responsible actors.

Fifth, it must be \emph{compatible with existing law}. A useful reform builds on company law, administrative law, data protection, and sector-specific supervision rather than attempting a total reconstruction.

Concretely, one can imagine a legal digital personality as a registration and attribution layer attached to certain high-impact AI deployments. It would not own rights in a general sense. It would not vote, contract freely, or enjoy constitutional standing. But it could serve as a legally recognised anchor for documentation, deployment history, designated responsibility points, and selected reporting or fiscal consequences. The decisive point is that the legal center of gravity remains with the accountable human or organisational actors.

\section{Objections and Limits}
The first objection is that the concept is unnecessary. Perhaps actor-based regulation already does enough. This objection is serious. In many sectors, it may indeed be correct. A limited legal digital personality should therefore be proposed only where ordinary categories prove persistently inadequate.

The second objection is conceptual inflation. Once one uses the language of personality, legal and public discourse may slide toward unwarranted anthropomorphism. That risk can be reduced only by rigid doctrinal boundaries and explicit statutory language denying any broader implication.

The third objection is strategic. Firms might use the category to shift blame onto a technical entity. For this reason, any such regime must be designed so that the digital personality supplements, rather than replaces, human and corporate obligations.

The fourth objection is comparative. Different jurisdictions may have little appetite for such innovation. This is especially true where policy already favors sector-specific rules and non-binding instruments. The concept should therefore be treated as a research hypothesis and not as an inevitable legislative endpoint.

\section{Conclusion}
The 2019 bachelor thesis made two moves. It rejected robot taxation in most ordinary settings, and it suggested that any serious attempt to tax or govern machine-generated activity would eventually encounter the problem of legal status. The first move remains persuasive. The second was insightful, but it now needs reinterpretation.

In 2026, the most defensible conclusion is not that AI should become a full legal person. Current legal developments in Europe and Switzerland point in another direction: they regulate providers, deployers, and risky uses while preserving human-centric responsibility. Yet the status problem has not vanished. It has narrowed. The remaining question is whether law sometimes needs a modest attribution interface to govern complex AI systems more precisely.

That is the value of legal digital personality. Properly understood, it is not a rival to actor-based regulation, nor a celebration of machine personhood. It is a limited legal technique for cases in which ordinary categories are no longer sufficiently precise. Read in this way, the 2019 thesis does not culminate in speculative futurism. It becomes the starting point for a disciplined 2026 research program about responsibility, attribution, and public control in the age of AI.

\bibliographystyle{plainnat}
\begin{thebibliography}{99}
\bibitem{thesis2019} Graz, David. \textit{Taxation des robots}. Bachelor thesis, 2019.
\bibitem{oberson2017} Oberson, Xavier. \textit{Taxing Robots? From the Emergence of an Electronic Ability to Pay to a Tax on Robots or the Use of Robots}. 2017.
\bibitem{floridi2017} Floridi, Luciano. \textit{Robots, Jobs, Taxes, and Responsibilities}. 2017.
\bibitem{acemoglu2018} Acemoglu, Daron and Restrepo, Pascual. \textit{Artificial Intelligence, Automation and Work}. 2018.
\bibitem{guerreiro2019} Guerreiro, Joao, Rebelo, Sergio and Teles, Pedro. \textit{Should Robots Be Taxed?} 2019.
\bibitem{thuemmel2018} Thuemmel, Uwe. \textit{Optimal Taxation of Robots}. 2018.
\bibitem{euai2026} European Commission materials on the AI Act, accessed 2026.
\bibitem{swissai2025} Swiss Federal Council materials on AI regulation, 2025-2026.
\bibitem{coe2024} Council of Europe Framework Convention on Artificial Intelligence, 2024.
\end{thebibliography}

\end{document}
